\documentclass[12pt]{article}
\usepackage[T2A]{fontenc}
\usepackage[utf8]{inputenc}
\usepackage[russian]{babel}
\usepackage{geometry}
\usepackage{booktabs}
\usepackage{longtable}
\usepackage{array}
\usepackage{graphicx}
\usepackage{float}
\geometry{a4paper, margin=2cm}

\begin{document}

\begin{center}
{\Large Отчет по лабораторной работе №1}\\
Исследование использования AI-инструментов в разработке программного обеспечения
\end{center}

\section*{Актуальность}
AI-инструменты быстро входят в повседневную практику разработки программного обеспечения, влияя на скорость работы, качество кода и требования к инженерным навыкам. Поэтому анализ частоты использования AI, уровня доверия к таким инструментам и степени их интеграции в процессы разработки представляет практический интерес.

\section*{Анкета и выборка}
Для исследования была составлена анкета из восьми вопросов с закрытыми вариантами ответов. Доверительная вероятность выбрана на уровне 95\%. При предположении максимальной дисперсии доли $p=0.5$ и допустимой ошибке 10\% минимальный объем выборки составляет 97 респондентов:
\[
n = \frac{z^2 p (1-p)}{E^2} = \frac{1.96^2 \cdot 0.5 \cdot 0.5}{0.1^2} \approx 97.
\]
В работе использована выборка из 100 респондентов.

\begin{longtable}{|m{1.5cm}|m{6cm}|m{7cm}|}
\hline
Код & Вопрос & Варианты ответов \\ \hline
Q1 & Как часто вы используете AI-инструменты в разработке? & 1 - Не использую; 2 - Редко, для отдельных задач; 3 - Регулярно несколько раз в неделю; 4 - Практически ежедневно \\ \hline
Q2 & Насколько глубоко AI встроен в ваш рабочий процесс? & 1 - Не встроен; 2 - Используется на одном этапе работы; 3 - Используется на нескольких этапах; 4 - Покрывает почти весь цикл разработки \\ \hline
Q3 & Насколько уверенно вы формулируете запросы к AI? & 1 - Нет опыта; 2 - Базовый уровень; 3 - Уверенный уровень; 4 - Продвинутый уровень \\ \hline
Q4 & Насколько вы доверяете AI-коду после собственной проверки? & 1 - Почти не доверяю; 2 - Доверяю только простым фрагментам; 3 - Доверяю типовым решениям; 4 - Часто принимаю AI-решение после ревью \\ \hline
Q5 & Как ваша команда относится к применению AI? & 1 - Использование запрещено; 2 - Отношение нейтральное; 3 - Использование разрешено и ограниченно регламентировано; 4 - Использование поощряется \\ \hline
Q6 & Насколько технически интегрированы AI-инструменты в вашу среду? & 1 - Интеграции нет; 2 - Использую только веб-интерфейсы; 3 - Есть плагины/расширения в IDE; 4 - Есть интеграция в IDE и/или командные процессы \\ \hline
Q7 & Какой прирост продуктивности вы ощущаете от AI? & 1 - Почти никакого; 2 - Небольшой, до 10%; 3 - Заметный, 10-30%; 4 - Существенный, более 30% \\ \hline
Q8 & Планируете ли вы расширять использование AI в ближайший год? & 1 - Нет; 2 - Скорее нет, чем да; 3 - Скорее да, чем нет; 4 - Да, планирую активное расширение \\ \hline
\end{longtable}

\section*{Доверительные интервалы для вопроса Q1}
Для долей ответов использован интервал Уилсона при доверительной вероятности 95\%.

\begin{table}[H]
\centering
\begin{tabular}{|c|m{5cm}|c|c|c|}
\hline
Код & Ответ & Кол-во & Доля & 95\% ДИ \\ \hline
1 & Не использую & 23 & 0.23 & [0.16; 0.32] \\ \hline
2 & Редко, для отдельных задач & 29 & 0.29 & [0.21; 0.39] \\ \hline
3 & Регулярно несколько раз в неделю & 19 & 0.19 & [0.13; 0.28] \\ \hline
4 & Практически ежедневно & 29 & 0.29 & [0.21; 0.39] \\ \hline
\end{tabular}
\caption{Доли ответов по вопросу Q1}
\end{table}

\section*{Подготовка данных и кластеризация}
Ответы были закодированы числами от 1 до 4, после чего выполнена нормировка по формуле min-max:
\[
x_{norm} = \frac{x - x_{min}}{x_{max} - x_{min}}.
\]
Для вычисления расстояния между анкетами использована евклидова метрика. Для объединения кластеров применен метод Уорда, минимизирующий прирост внутрикластерной дисперсии. По дендрограмме выбран уровень отсечения $h \approx 3.618$, что соответствует выделению трех кластеров.

\section*{Профили кластеров}
В таблице приведены средние значения ответов по вопросам для каждого кластера.

\begin{table}[H]
\centering
\small
\begin{tabular}{|c|c|c|c|c|c|c|c|c|c|}
\hline
Кластер & Размер & Q1 & Q2 & Q3 & Q4 & Q5 & Q6 & Q7 & Q8 \\ \hline
1 & 36 & 1.47 & 1.33 & 1.78 & 1.86 & 2.00 & 1.78 & 1.69 & 2.00 \\ \hline
2 & 25 & 3.80 & 3.68 & 3.64 & 3.68 & 3.52 & 3.68 & 3.68 & 3.88 \\ \hline
3 & 39 & 2.72 & 3.08 & 2.67 & 2.79 & 3.05 & 2.82 & 2.51 & 3.05 \\ \hline
\end{tabular}
\caption{Средние значения ответов по кластерам}
\end{table}

\noindent Интерпретация кластеров:
\begin{itemize}
\item Кластер 1 (размер 36): Кластер консервативных пользователей: AI применяется редко, интеграция слабая, ожидаемый эффект умеренный или низкий.
\item Кластер 2 (размер 25): Кластер активных пользователей AI: инструменты применяются почти ежедневно, глубоко встроены в процесс и воспринимаются как источник высокого прироста продуктивности.
\item Кластер 3 (размер 39): Кластер прагматичных пользователей: AI уже используется регулярно, но без полной зависимости; респонденты видят практическую пользу и готовы расширять использование.
\end{itemize}

\section*{Дендрограмма}
Дендрограмма строится в Wolfram Mathematica на основе нормированных данных. Файл с кодом: \texttt{lab1\_dendrogram.wl}. При выполнении кода формируется рисунок \texttt{dendrogram.png} с подписями осей и линией уровня отсечения.

\section*{Выводы}
Проведенный анализ показывает наличие трех устойчивых групп респондентов: консервативных, прагматичных и продвинутых пользователей AI. Наиболее информативными для разделения кластеров оказались частота использования AI, глубина интеграции в рабочий процесс, техническая интеграция инструментов и ожидаемый прирост продуктивности. Полученная сегментация может использоваться для настройки корпоративной политики внедрения AI и выбора образовательных программ для разработчиков.

\end{document}
